\documentclass[11pt]{article}

\usepackage{extsizes}

\usepackage[a4paper, margin=1in]{geometry}
\usepackage{hyperref}
\usepackage[style=ieee]{biblatex}
\usepackage{listings}
\usepackage{xcolor}

\definecolor{codegreen}{rgb}{0,0.6,0}
\definecolor{codegray}{rgb}{0.5,0.5,0.5}
\definecolor{codepurple}{rgb}{0.58,0,0.82}
\definecolor{backcolour}{rgb}{0.95,0.95,0.92}

\lstdefinestyle{mystyle}{
    backgroundcolor=\color{backcolour},
    commentstyle=\color{codegreen},
    keywordstyle=\color{magenta},
    numberstyle=\tiny\color{codegray},
    stringstyle=\color{codepurple},
    basicstyle=\ttfamily\footnotesize,
    breakatwhitespace=false,
    breaklines=true,
    captionpos=b,
    keepspaces=true,
    numbers=left,
    numbersep=5pt,
    showspaces=false,
    showstringspaces=false,
    showtabs=false,
    tabsize=2
}

\lstset{style=mystyle}

\usepackage{parskip}

\title{SeaSTAR Documentation}
\author{A Baldwin}
\date{February 2026}

\begin{document}

\maketitle

\begin{abstract}

\end{abstract}

\tableofcontents

\newpage
\section{Installing SeaSTAR}

The easiest way to install SeaSTAR is via PyPI into a Python Virtual Environment:

\begin{lstlisting}
[alewin@nike ~]$ python3 -m venv seastarvenv
[alewin@nike ~]$ source ~/seastarvenv/bin/activate
(seastarvenv) [alewin@nike ~]$ pip install seastartool
\end{lstlisting}

Whenever you are using SeaSTAR like this, you need to first enable the Virtual Environment with the \texttt{source \~{}/seastarvenv/bin/activate} command. You should now be able to run the SeaSTAR CLI as so:

\begin{lstlisting}
(general_venv) [alewin@nike docs]$ seastar

SeaSTAR
Sea-faring System for Tagging, Attribution and Redistribution

Copyright 2025, A Baldwin <alewin@noc.ac.uk>, National Oceanography Centre
This program comes with ABSOLUTELY NO WARRANTY. This is free software,
and you are welcome to redistribute it under the conditions of the
GPL version 3 license.

ERROR
No command specified

Common usage:

seastar --gui                           Launch the SeaSTAR GUI
seastar <job_name> [flags]              General structure

With any SeaSTAR job you can use --logfile <file_path> to send all log output
to a file in addition to the console. This log output will also include the date
and time, along with all arguments used to invoke SeaSTAR.

Use "seastar <job_name> --help" to get specific help for a given job.
The following jobs are avaliable:

ecotaxa_to_crab                         EcoTaxa to CRAB compatible upload

    Creates CRAB compatible parquet files for uploading images or updating
    annotations in CRAB.

ifcb_v4_features                        IFCB feature extractor (WHOI v4 compatible)

    Extracts morphological feature information from raw IFCB data. Is
    compatible with outputs of the V4 WHOI classifier.

ifcb_to_ecotaxa                         EcoTaxa IFCB Data Import

    Creates EcoTaxa TSV or ZIP files for uploading images or updating metadata
    from raw IFCB data, and/or morphological feature files.

\end{lstlisting}

\section{Extracting features from IFCB data}

\begin{lstlisting}
\end{lstlisting}

\end{document}
